\section{Numerical Experiments}
% ============================================================

\subsection{Verification Against the Exact Matrix}
We integrate the prescribed $X_0$ over $[0,2]$ on fixed uniform grids
and measure the full-matrix terminal error
\[
E_h(2)=\|X_h(2)-X_{\rm exact}(2)\|_F.
\]
The limiting orthogonal factor is not the finite-time reference.
Euler and Implicit Euler use $N\in\{40,80,160,320\}$.
RK4 also uses $N\in\{640,1280,2560\}$.
The finest three RK4 grids are used for the displayed slope because
the coarse-grid ratios are affected by the rapid initial transient.
The smallest measured error is near floating-point round-off scale,
so the fitted fine-grid slope is a descriptive validation diagnostic;
no flat round-off floor has been observed on these grids. Table~\ref{tab:order_results} and
Fig.~\ref{fig:convergence} report the measured orders.

\begin{table}[htbp]
\centering
\begin{tabular}{lrrr}
\toprule
Method & Fitted slope & Finest $h$ & Finest terminal error \\
\midrule
Explicit Euler & $1.010$ & $0.00625$ & $4.6141\times10^{-4}$ \\
RK4 & $3.927$ & $0.00078125$ & $7.3968\times10^{-14}$ \\
Implicit Euler & $1.014$ & $0.00625$ & $4.6429\times10^{-4}$ \\
\bottomrule
\end{tabular}
\caption{Finite-time full-matrix Frobenius error. The complete
grid-by-grid values and adjacent observed orders are in
\texttt{data/convergence.csv}.}
\label{tab:order_results}
\end{table}

\begin{figure}[htbp]
\centering
\includegraphics[width=0.76\textwidth]{convergence.png}
\caption{Terminal full-matrix error versus dimensionless step size, with
$h^1$ or $h^4$ reference slopes and fitted slopes from the refined grids.
The shaded 95\% regression bands and reported slope standard errors are
descriptive least-squares diagnostics for deterministic mesh data, not
probabilistic uncertainty in the ODE solution. The agreement verifies the
predicted global orders against the finite-time exact matrix.}
\label{fig:convergence}
\end{figure}

A separate \texttt{solve\_ivp} Radau run with relative tolerance
$10^{-12}$ and absolute tolerance $10^{-14}$ differs from the exact
SVD matrix at $t=2$ by $1.16\times10^{-13}$ in Frobenius norm. This
cross-check and the analytic formula are independent of the three
implemented time-stepping methods.

\subsection{Stability Experiments}
The initial frozen-Jacobian estimates
$h_{\rm crit}^{\rm EE}\approx0.0769$ and
$h_{\rm crit}^{\rm RK4}\approx0.1071$ predict the onset of possible
amplification in the fastest contracting direction. We test selected
fixed steps $h=2/N$ with $N=6,7,8,10,12,16,20,24,32,40,50,64,80$
and perturb the initial state by $10^{-7}$ in that direction.
At $h=0.10$, one Euler step amplifies this perturbation by about $1.60$,
although the full numerical path still finishes. At $h=2/7$, Euler
increases $\Phi$ by $31.9$ on its first step and the run subsequently
fails. RK4 at $h=1/3$ also fails after a large first-step energy rise,
whereas Implicit Euler finishes that tested grid. These observations
support the predicted local risk but do not turn the frozen-eigenvalue
cutoffs into nonlinear global thresholds.

\begin{figure}[htbp]
\centering
\includegraphics[width=0.85\textwidth]{stability_sweep.png}
\caption{Left: one-step amplification of a fast-direction initial
perturbation. Right: terminal error for completed runs on log--log axes. The labelled
vertical lines show Euler and RK4 limits from the initial frozen
Jacobian; they are local estimates, so the nonlinear sweep is needed
to assess their relevance. Energy changes and failed-run statuses are in
\texttt{stability\_sweep.csv}.}
\label{fig:stability_sweep}
\end{figure}

\subsection{Adaptivity}
For all three base methods, the step-doubling trials start at $h=0.15$
with $\mathrm{atol}=10^{-8}$ and $\mathrm{rtol}=10^{-6}$.
The fine result is accepted only when its normalized estimated local
error is at most one. Table~\ref{tab:adaptive_results} and
Fig.~\ref{fig:adaptive_steps} show that the accepted step size genuinely
changes. A very short final step may be caused by clipping to reach
$t=2$ exactly. These local tolerances do not provide a strict bound on
the terminal global error, which we measure separately against the
exact SVD matrix.

\begin{table}[htbp]
\centering
\begin{tabular}{lrrrr}
\toprule
Method & Accepted & Rejected & Max.\ accepted $E$ & Terminal error \\
\midrule
Explicit Euler & 1680 & 5 & $0.842$ & $5.08\times10^{-5}$ \\
RK4 & 17 & 2 & $0.754$ & $6.89\times10^{-7}$ \\
Implicit Euler & 1678 & 5 & $0.843$ & $5.09\times10^{-5}$ \\
\bottomrule
\end{tabular}
\caption{Step-doubling runs with the same local tolerance settings.
The numbers of accepted steps alone are not a work comparison:
each trial computes one full and two half steps, and rejected trials
also consume evaluations.}
\label{tab:adaptive_results}
\end{table}

\begin{figure}[htbp]
\centering
\includegraphics[width=0.78\textwidth]{adaptive_steps.png}
\caption{Accepted step length versus dimensionless time. Step doubling
reduces the step during the rapid initial transient; the accepted-step
data and estimated local error are in \texttt{adaptive\_steps.csv}.}
\label{fig:adaptive_steps}
\end{figure}

\subsection{Main Results and Comparison of Methods}
\paragraph{Geometry and Lyapunov diagnostics.}
On the RK4 trajectory with $h=0.005$, every sampled singular value
moves toward one without crossing it, and $\Phi(X)$ decreases at every
stored step. The numerical orthogonality defect at $t=2$ is
$0.305932824226$, compared with the nonzero exact finite-time value
$0.305932824210$. Figure~\ref{fig:trajectory_diagnostics} also shows
the finite-difference check of
$\mathrm d\Phi/\mathrm dt=-\|X(I-X^TX)\|_F^2$.
Its maximum absolute residual is about $0.240$, concentrated in the
initial fast transient; the maximum residual scaled by
$\max(1,|\mathrm d\Phi/\mathrm dt|)$ is $4.69\times10^{-4}$.
This is a discrete diagnostic, not a direct test that the continuous
identity fails.

\begin{figure}[htbp]
\centering
\includegraphics[width=\textwidth]{trajectory_diagnostics.png}
\caption{Singular values, Lyapunov energy, finite-time orthogonality
defect, and the discrete identity residual along a resolved RK4
trajectory. Dashed curves give the finite-time exact SVD reference for
the first three panels; the fourth panel explicitly labels the
midpoint Lyapunov residual. Their overlap confirms accuracy while the
nonzero exact defect explains why finite-time orthogonality is not expected.}
\label{fig:trajectory_diagnostics}
\end{figure}

\paragraph{Cost at comparable accuracy.}
We select fixed grids that yield terminal full-matrix errors within
$5.4\%$ of one another, then count RHS evaluations as the common
work measure. Table~\ref{tab:cost_accuracy} and
Fig.~\ref{fig:cost_accuracy} show that RK4 needs fewer RHS calls for
this target, while Implicit Euler incurs many nonlinear-solver calls.
Jacobian evaluations and Newton updates are reported separately;
the counts do not include matrix factorization, Python overhead, or
plotting and should not be interpreted as wall-clock timings.

\begin{table}[htbp]
\centering
\begin{tabular}{lrrrr}
\toprule
Method & $N$ & Terminal error & RHS calls & Jacobian/Newton updates \\
\midrule
Explicit Euler & 160 & $9.2066\times10^{-4}$ & 160 & $0/0$ \\
RK4 & 17 & $8.8502\times10^{-4}$ & 68 & $0/0$ \\
Implicit Euler & 160 & $9.3194\times10^{-4}$ & 862 & $351/351$ \\
\bottomrule
\end{tabular}
\caption{Counted work at closely matched terminal error.}
\label{tab:cost_accuracy}
\end{table}

\begin{figure}[htbp]
\centering
\includegraphics[width=0.86\textwidth]{cost_accuracy.png}
\caption{Achieved terminal full-matrix error versus counted RHS calls
over several fixed grids. Stars mark the three table entries with errors
near $9\times10^{-4}$; the grey band spans their actual errors. RK4
uses 68 RHS calls at this target, versus 160 and 862 for Explicit and
Implicit Euler. Jacobian and Newton costs are reported separately in
\texttt{cost\_accuracy.csv}, so RHS count is not wall-clock time.}
\label{fig:cost_accuracy}
\end{figure}

\paragraph{Rank-deficient initial state.}
Replacing the first initial singular value by zero and integrating
over $[0,8]$ with RK4 at $h=0.005$ leaves the smallest computed
singular value below $7.1\times10^{-15}$ throughout the stored path.
The two nonzero terminal singular values are approximately
$1.0000000500$ and $1.0000000276$; the final error against the
finite-time exact matrix is $3.86\times10^{-14}$.
The final orthogonality defect is approximately one, consistent with
a rank-two partial isometry rather than an orthogonal limit
(Fig.~\ref{fig:rank_deficient}).

\begin{figure}[htbp]
\centering
\includegraphics[width=0.78\textwidth]{rank_deficient.png}
\caption{The zero singular mode remains at round-off scale while the
other modes approach one, agreeing with the dashed exact SVD curves.
This distinguishes the partial-isometry limit from an orthogonal one.
Values are in
\texttt{rank\_deficient.csv}.}
\label{fig:rank_deficient}
\end{figure}


\subsection{Verification Coverage}
The runnable checks in \path{test/} comprise eleven test functions.
They check the model and full implicit-residual Jacobians by centred
finite differences, both complete eigenvalue spectra, the exact matrix
against tight-tolerance Radau, and a diagonal scalar sanity case.
All three methods are checked for fixed-grid order, genuinely varying
adaptive steps and accepted normalized estimated local errors at most
one. Their well-resolved trajectories are also checked for singular
values approaching one without crossing and non-increasing Lyapunov
energy. The rank-deficient check verifies the zero mode and convergence
toward the rank-two partial isometry. Independent subprocesses compare
both language versions for all three fixed and adaptive solvers.
The checks pass on the recorded environment.

The RK4 order test uses $N=640$ and $1280$, avoiding the finest
experiment grid whose error is nearer round-off scale. Passing these
checks supports the tested configurations; it does not turn the
local adaptive estimator into a universal global error bound.
The full numerical evidence and CSV interpretation are supplied in
Appendix~\ref{sec:appendix_c}.

\subsection{Limitations and Robustness}
The initial $h\lambda$ overlay is a frozen-Jacobian warning, not a
sufficient global nonlinear stability condition. An accepted
step-doubling estimate is likewise a local indicator rather than a
proof of global accuracy or stability. Large steps can destroy
monotonicity of $\Phi$ or cause an explicit run to fail; Newton may also
fail to find a decreasing candidate. The fixed-grid solver reports
such failures instead of changing the grid used to estimate order,
whereas the adaptive solver retries from the last accepted state.
Our matched-accuracy cost study counts RHS calls and reports Newton
work separately; a timing study would be needed to compare actual
runtime on a specified machine.

% ============================================================
